
% This must be in the first 5 lines to tell arXiv to use pdfLaTeX, which is strongly recommended.
\pdfoutput=1
% In particular, the hyperref package requires pdfLaTeX in order to break URLs across lines.

\documentclass[11pt]{article}

% Change "review" to "final" to generate the final (sometimes called camera-ready) version.
% Change to "preprint" to generate a non-anonymous version with page numbers.
\usepackage[review]{acl}

% Standard package includes
\usepackage{times}
\usepackage{latexsym}
\usepackage[T1]{fontenc}
\usepackage[utf8]{inputenc}
\usepackage{microtype}
\usepackage{inconsolata}

% ── Extra packages ────────────────────────────────────────────
\usepackage{graphicx}
\usepackage{booktabs}
\usepackage{multirow}
\usepackage{amsmath,amssymb}
\usepackage{hyperref}
\usepackage{xcolor}
\usepackage{enumitem}

\title{Preserving Epistemic Uncertainty: Token-Level Masking Prevents Self-Distillation Degradation in Mathematical Reasoning}

% Anonymous submission — authors hidden
\author{}

\begin{document}
\maketitle

\begin{abstract}
Self-distillation has emerged as an effective LLM post-training paradigm, yet recent work reveals a critical failure mode: when applied to mathematical reasoning, it degrades out-of-distribution performance by suppressing \textit{epistemic verbalization}—the model's expression of uncertainty through tokens such as ``Wait'' and ``Hmm.''
We provide the first mechanistic diagnosis of this phenomenon, showing via token-level probability probes that a teacher conditioned on reference solutions systematically undervalues epistemic tokens, assigning them less than half the probability of non-epistemic tokens.
This concentrated KL divergence forces the student to abandon uncertainty expression, collapsing its reasoning capability.
We propose a minimal fix: token-level masking of epistemic markers during distillation, which excludes uncertainty positions from the KL loss at zero overhead.
Experiments across two model scales and two reference-solution types demonstrate that masking reduces AIME24 degradation from $-20\%$ to $-8\%$ and completely preserves AIME25 accuracy.
Our diagnostic and interventional evidence converges on a clear message: protecting epistemic verbalization is both necessary and sufficient to prevent self-distillation from degrading mathematical reasoning.
\end{abstract}

% ═══════════════════════════════════════════════════════════════
\section{Introduction}
% ═══════════════════════════════════════════════════════════════

Self-distillation \cite{snell2022learning,agarwal2024onpolicy,zhao2026self,hubotter2026reinforcement} has recently gained traction as a post-training paradigm for LLMs.
In this framework, the same model acts as both student and teacher under asymmetric conditioning: the teacher is provided with ground-truth solutions or environment feedback as privileged information, while the student receives only the problem.
Training minimizes the divergence between the student's and teacher's next-token distributions, effectively distilling the knowledge present in the teacher's privileged context into the student's parameters.

Two prominent instantiations have emerged.
\textbf{On-Policy Self-Distillation} (OPSD; \cite{zhao2026self}) conditions the teacher on a verified reference solution from a reasoning dataset, using the same model to rationalize the solution and guide the student's own rollouts.
\textbf{Self-Distillation Policy Optimization} (SDPO; \cite{hubotter2026reinforcement}) provides the teacher with rich environment feedback (e.g., runtime errors, test failures) and uses successful rollouts from the same batch as privileged context for failed attempts.

Both methods report substantial gains in their target domains.
OPSD improves AIME24 from 51.5\% to 57.2\% on Qwen3-1.7B \cite{zhao2026self}.
SDPO achieves 48.8\% on LiveCodeBench v6 with Qwen3-8B, outperforming GRPO by a wide margin with 4$\times$ fewer generations \cite{hubotter2026reinforcement}.
These results suggest that self-distillation is a broadly applicable and powerful post-training strategy.

However, this paradigm harbors a striking failure mode.
When applied to mathematical reasoning, self-distillation can severely degrade out-of-distribution performance.
We trace the root cause to the suppression of \textit{epistemic verbalization}: the model's explicit expression of uncertainty through tokens like ``Wait,'' ``Hmm,'' ``perhaps,'' and ``maybe.''
When the teacher is conditioned on the correct solution, it produces confident, concise reasoning with minimal expressed uncertainty.
The student learns to imitate this style, losing the ability to explore and recover from intermediate errors.

This apparent contradiction—self-distillation succeeds on code yet fails on math—raises a fundamental question:
\textit{Why does self-distillation degrade reasoning in some domains but not others?}

In this paper, we provide a systematic diagnosis and a targeted solution.
We show that the degradation is driven by a concentrated, measurable mechanism: the teacher's privileged context causes it to systematically underestimate the probability of epistemic tokens, creating a per-token KL divergence gap that forces the student to suppress uncertainty.
Critically, we demonstrate that this mechanism operates with varying intensity depending on the \textit{reference context}: concise reference solutions produce mild suppression, while full chain-of-thought traces (themselves containing epistemic markers) create extreme suppression that polarizes the student's behavior.

Armed with this mechanistic understanding, we propose a minimal, zero-overhead intervention: \textit{token-level masking of epistemic markers} during the self-distillation loss computation.
By excluding positions where the student generates an epistemic token from the KL divergence, we prevent the teacher from forcing an overconfident reasoning style while preserving normal distillation everywhere else.

Our contributions are:
\begin{enumerate}[leftmargin=*,nosep]
    \item \textbf{Mechanistic diagnosis}: We conduct token-level probability probes revealing that the teacher assigns 2.2$\times$ lower probability to epistemic tokens (0.34 vs.\ 0.75), creating a concentrated KL gap that drives reasoning degradation.
    \item \textbf{Domain and context analysis}: We show that reference context type (concise vs.\ full trace) modulates degradation severity, explaining why self-distillation succeeds in code (no epistemic tokens needed) but fails in math reasoning.
    \item \textbf{Epistemic token masking}: A simple, training-free fix that preserves epistemic tokens during distillation, reducing AIME24 degradation from $-20\%$ to $-8\%$ and completely eliminating AIME25 degradation.
    \item \textbf{Multi-scale validation}: Experiments across two model scales (1.5B and 7B) and two reference-solution types demonstrate consistent benefits.
\end{enumerate}


% ═══════════════════════════════════════════════════════════════
\section{Background and Related Work}
% ═══════════════════════════════════════════════════════════════

\subsection{Self-Distillation for LLM Post-Training}

Self-distillation \cite{snell2022learning} uses a single model as both teacher and student under different conditioning contexts.
Given an input $x$ and ground-truth answer $y^*$, the student generates a sequence $\hat{y} \sim \pi_\theta(\cdot \mid x)$.
The teacher policy $\pi_\theta^T$ is obtained by conditioning the same model on privileged information $s$ (e.g., a reference solution or environment feedback):
$\pi_\theta^T(\cdot \mid x, s) = \pi_\theta(\cdot \mid x, s)$.
Training minimizes the KL divergence between student and teacher next-token distributions:

\begin{equation}
{\small
    \mathcal{L}_{\text{SD}}(\theta) = \sum_{t} \text{KL}\bigl(
        \pi_\theta(\cdot \mid x, \hat{y}_{<t}) \;\|\;
        \text{sg}(\pi_\theta(\cdot \mid x, s, \hat{y}_{<t}))
    \bigr)
}
\end{equation}

where $\text{sg}(\cdot)$ denotes stop-gradient, preventing the teacher from being updated.
Recent instantiations include OPSD \cite{zhao2026self}, which uses fixed reference solutions as teacher context, and SDPO \cite{hubotter2026reinforcement}, which uses dynamic environment feedback and successful in-batch rollouts.

\subsection{Epistemic Verbalization and Reasoning}

Epistemic verbalization refers to the explicit externalization of uncertainty during reasoning through tokens such as ``wait,'' ``hmm,'' ``perhaps,'' and ``maybe.''
Strong reasoning models like DeepSeek-R1 frequently use these tokens, and their removal degrades reasoning performance.
When applied to self-distillation, we observe that the teacher's privileged access to the correct solution causes it to reduce epistemic markers substantially.
When the student is trained to imitate these confident traces, it loses the ability to express and act on uncertainty, leading to catastrophic out-of-distribution degradation.
Crucially, we show that the per-token KL divergence between student and teacher on epistemic tokens is concentrated and disproportionate compared to the corpus-wide mean.

\subsection{Reasoning Compression}

Several concurrent works aim to shorten reasoning traces while preserving accuracy \cite{shrivastava2025sample,sang2026onpolicy,deng2026conpress,yang2026accordion,luo2026compress}.
Our work complements these efforts: rather than designing new compression algorithms, we identify \textit{which} tokens should be protected from compression and provide a simple mechanism to do so within any self-distillation pipeline.


% ═══════════════════════════════════════════════════════════════
\section{Diagnosis: Why Self-Distillation Degrades Reasoning}
% ═══════════════════════════════════════════════════════════════

Before proposing a solution, we conduct a systematic diagnosis of the degradation mechanism.
Our analysis proceeds in three stages: (1) a domain-level contrast showing that self-distillation degradation is specific to reasoning tasks, (2) a token-level probability probe revealing \textit{how} the teacher suppresses epistemic tokens, and (3) a reference-context ablation demonstrating that degradation severity depends on the type of privileged information.

\subsection{Domain Contrast: Code Succeeds, Math Fails}

Table~\ref{tab:domain_contrast} summarizes the performance of self-distillation across domains.
In code generation (LiveCodeBench v6), SDPO achieves 48.8\% accuracy, substantially outperforming GRPO (41.2\%) with $4\times$ fewer generations \cite{hubotter2026reinforcement}.
In scientific reasoning and tool use, SDPO similarly improves over RLVR baselines \cite{hubotter2026reinforcement}.
However, in mathematical reasoning (AIME), OPSD produces only modest gains (51.5\% $\to$ 57.2\% on Qwen3-1.7B) while SDPO causes catastrophic degradation (55\% $\to$ $<$20\% on DeepSeek-7B).

\begin{table*}[t]
\centering
\begin{tabular}{lccc}
\toprule
\textbf{Domain} & \textbf{Method} & \textbf{Performance} & \textbf{Effect} \\
\midrule
Code (LiveCodeBench) & SDPO & 48.8\% (vs.\ GRPO 41.2\%) & $\uparrow$ Strong gain \\
Science \& Tools & SDPO & 70.2\% (vs.\ GRPO 66.6\%) & $\uparrow$ Moderate gain \\
Math (AIME, OPSD) & OPSD & 51.5\% $\to$ 57.2\% & $\uparrow$ Mild gain \\
Math (AIME, SDPO) & SDPO & 55\% $\to$ $<$20\% & $\downarrow$\textbf{ Severe degradation} \\
\bottomrule
\end{tabular}
\caption{Self-distillation performance across domains. Code and science tasks benefit; mathematical reasoning degrades.}
\label{tab:domain_contrast}
\end{table*}

Why does this domain split exist?
We hypothesize that code generation and scientific tool use rely primarily on syntactic patterns, memorized APIs, and structured reasoning, where epistemic uncertainty markers are unnecessary.
In contrast, mathematical reasoning requires extended chains of inference, self-verification, and exploration of alternative paths—processes that are explicitly mediated by epistemic verbalization.
When the teacher suppresses these tokens, the student loses the cognitive scaffolding needed for complex reasoning.

\subsection{Token-Level Probability Probe}

To quantify the teacher's suppression of epistemic tokens, we conduct a controlled forward-pass experiment.
Using DeepSeek-R1-Distill-Qwen-1.5B, we sample one student-generated response on an AIME-style problem (1024 tokens), then compute both the teacher (conditioned on a COT reference solution) and student (conditioned only on the problem) next-token probability distributions at every position.

Table~\ref{tab:prob_probe} reports the mean probability assigned to epistemic vs.\ non-epistemic tokens.

\begin{table}[h]
\centering
\small
\begin{tabular}{lccc}
\toprule
\textbf{Condition} & \textbf{Epistemic} & \textbf{Non-Epi.} & \textbf{Ratio} \\
\midrule
Teacher (COT) & 0.336 & 0.751 & 0.45$\times$ \\
Student (plain)    & 0.417 & 0.761 & 0.55$\times$ \\
\bottomrule
\end{tabular}
\caption{Mean next-token probability for epistemic vs.\ non-epistemic tokens. The teacher assigns only 45\% as much probability to epistemic tokens.}
\label{tab:prob_probe}
\end{table}

The teacher assigns only \textbf{0.336} mean probability to epistemic tokens—2.2$\times$ lower than the 0.751 assigned to non-epistemic tokens.
The student, in contrast, maintains a higher epistemic probability of 0.417.
This creates a concentrated KL divergence gap at epistemic token positions, where the student naturally expresses uncertainty that the teacher deems unlikely.

\subsection{Top-10 Token Analysis}

To further characterize the teacher's behavior, we examine the top-10 most probable next-token predictions at each epistemic position.
Table~\ref{tab:top10} shows representative examples.

\begin{table*}[t]
\centering
\small
\begin{tabular}{p{1.5cm}p{6cm}p{6cm}}
\toprule
\textbf{Pos.} & \textbf{Teacher Top-3 (probability)} & \textbf{Student Top-3 (probability)} \\
\midrule
\#41 ``Hmm'' & ``I'' (0.629), ``Hmm'' (0.232), ``It'' (0.066) & ``Hmm'' (0.453), ``I'' (0.453), ``It'' (0.069) \\
\#305 ``Hmm'' & ``So'' (0.361), ``Each'' (0.281), ``But'' (0.091) & ``So'' (0.410), ``But'' (0.118), ``Let'' (0.091) \\
\#261 ``Wait'' & ``Wait'' (0.637), ``Each'' (0.301), ``Hmm'' (0.022) & ``Wait'' (0.586), ``Each'' (0.244), ``Hmm'' (0.090) \\
\bottomrule
\end{tabular}
\caption{Top-3 token predictions by teacher and student at selected epistemic positions.}
\label{tab:top10}
\end{table*}

At position \#41, where the student naturally considers both ``Hmm'' and ``I'' equally (0.453 each), the teacher strongly prefers the non-epistemic ``I'' (0.629) and depresses ``Hmm'' to 0.232.
At position \#305, the teacher's top-10 predictions do not even include ``Hmm,'' instead preferring connective tokens like ``So,'' ``Each,'' and ``But.''
At position \#261, while both models rank ``Wait'' first, the student maintains ``Hmm'' as a viable alternative (0.090) while the teacher nearly eliminates it (0.022).

These examples illustrate a consistent pattern: \textbf{the teacher systematically prefers non-epistemic alternatives at positions where the student would naturally express uncertainty}.

\subsection{Reference Context Matters}

Finally, we investigate how the \textit{type} of reference solution affects epistemic suppression.
We compare four prompt conditions using the same 30-question probe, each with 3 samples:

\begin{itemize}[leftmargin=*,nosep]
    \item \textbf{A (Student)}: Plain problem, no privileged information.
    \item \textbf{B (GT Only)}: Problem + correct answer only.
    \item \textbf{C (Clean Solution)}: Problem + concise textbook solution.
    \item \textbf{D (Full COT)}: Problem + complete chain-of-thought reasoning trace.
\end{itemize}

\begin{table}[h]
\centering
\small
\begin{tabular}{lcc}
\toprule
\textbf{Condition} & \textbf{Mean Epi. Tokens} & \textbf{Acc.} \\
\midrule
A: Student (no context) & 17.17 & 33.3\% \\
B: GT only & 16.97 & 33.3\% \\
C: Clean reference solution & 10.77 & 56.7\% \\
D: Full COT trace & 3.50 & 90.0\% \\
\bottomrule
\end{tabular}
\caption{Epistemic token frequency and answer accuracy across prompt conditions (30 problems $\times$ 3 samples). Richer reference context suppresses epistemic tokens but inflates accuracy, masking reasoning degradation.}
\label{tab:ref_context}
\end{table}

Two conclusions emerge from Table~\ref{tab:ref_context}.
First, \textbf{clean reference solutions produce mild epistemic suppression} (37\% reduction, from 17.17 to 10.77).
Second, \textbf{full COT traces cause extreme suppression} (80\% reduction, to 3.50 tokens) but simultaneously inflate accuracy to 90\%—the model simply echoes the provided reasoning rather than producing its own.
This latter effect is pernicious: high in-distribution accuracy conceals the loss of independent reasoning capability.

Together, these diagnostic findings establish that (1) the teacher's token-level suppression of epistemic markers is a measurable, robust phenomenon, (2) the suppression is concentrated on a small, identifiable set of tokens, and (3) the severity depends on the richness of the reference context.
This mechanistic understanding directly motivates our intervention.

% ═══════════════════════════════════════════════════════════════
\section{Method: Epistemic Token Masking}
% ═══════════════════════════════════════════════════════════════

Our diagnosis reveals that self-distillation degradation is driven by a concentrated mechanism: the teacher's systematic suppression of a small, identifiable set of epistemic tokens.
This suggests a targeted fix: simply exclude these token positions from the distillation loss.

\subsection{Epistemic Token Set}

We define the set of epistemic markers as:
\begin{equation}
\begin{aligned}
    \mathcal{T} = \{&\text{wait}, \text{hmm}, \text{perhaps}, \text{maybe}, \text{actually}, \\
                    &\text{alternatively}, \text{seems}, \text{might}, \text{likely}, \text{check}\}
\end{aligned}
\end{equation}

To handle surface-form variation (capitalization, punctuation, surrounding whitespace), we enumerate common variants of each token and encode them with the model's tokenizer.
If a variant encodes to multiple token IDs, we keep only the last (the preceding tokens are typically whitespace artifacts).
This yields approximately 100--200 unique token IDs per model, cached and reused across training steps.

\subsection{Masked Self-Distillation Loss}

Let $\hat{y}_t$ be the student's generated token at position $t$, and let $\mathcal{I}$ be the set of token IDs corresponding to epistemic markers.
We define a binary token mask $M \in \{0,1\}^{T}$:
\begin{equation}
    M_t = \begin{cases}
        0 & \text{if } \hat{y}_t \in \mathcal{I} \\
        1 & \text{otherwise}
    \end{cases}
\end{equation}

The masked self-distillation loss becomes:
\begin{equation}
{\small
    \mathcal{L}_{\text{masked}}(\theta) = \frac{\sum_t M_t \cdot \ell_t}{\sum_t M_t},
    \quad
    \ell_t = \text{KL} \bigl( \pi_\theta^S \;\|\; \pi_\theta^T \bigr)
}
\label{eq:masked_loss}
\end{equation}

The denominator normalizes by the number of \textit{non-masked} tokens, ensuring that masking does not artificially reduce the loss magnitude.

\subsection{Implementation}

The masking is implemented as a lightweight, on-the-fly modification to the chunked KL divergence computation.
For each micro-batch, after computing the per-token KL divergence, we multiply by the pre-computed token mask before summing.
The mask is constructed from the response token IDs, which are already available from the forward pass.
No additional forward passes or memory are required, making the method a zero-overhead addition to any self-distillation pipeline.


% ═══════════════════════════════════════════════════════════════
\section{Experiments}
% ═══════════════════════════════════════════════════════════════

We evaluate epistemic token masking across two reference-solution types and two model scales to test both the generality of the degradation phenomenon and the effectiveness of our intervention.

\subsection{Setup}

\textbf{Models.}
We use DeepSeek-R1-Distill-Qwen at two scales: 1.5B and 7B parameters.
The 1.5B model is used for COT-reference experiments where sequence lengths are longer; the 7B model is used for clean-solution experiments.

\textbf{Data and Training.}
All experiments use the Openthoughts Math 30K dataset \cite{guha2025openthoughts}, which provides problems, reference solutions, and ground-truth answers.
We adopt the on-policy self-distillation (OPSD) configuration \cite{zhao2026self}: each training step samples 12 problems, generates 1 student rollout per problem, and applies self-distillation using a frozen reference teacher.
The teacher prompt includes a verified reference solution as privileged context (either the concise ``solution'' field or the full ``COT\_Reason'' chain-of-thought trace).
All experiments use forward KL divergence with a learning rate of $5 \times 10^{-6}$, micro-batch size of 1 per GPU, and train for 100 steps.
Experiments run on 8 NVIDIA A800-80GB GPUs.

\textbf{Evaluation.}
We evaluate on AIME 2024, AIME 2025, and MATH-500.
For AIME benchmarks, we sample 12 responses per problem (temperature 1.0, top-p 0.95) and report accuracy@12.
For MATH-500, we use greedy decoding (temperature 0) and report pass@1.
Evaluation occurs every 10 training steps.

\textbf{Configurations.}
We compare two training variants:
\begin{itemize}[leftmargin=*,nosep]
    \item \textbf{Naive OPSD}: Standard self-distillation without masking (sd\_mask\_mode=none).
    \item \textbf{Masked OPSD}: Our proposed method with token-identity masking (sd\_mask\_mode=token\_identity).
\end{itemize}

\subsection{COT Reference Solution Experiments (1.5B)}

We first evaluate the most challenging setting: using full chain-of-thought reasoning traces as the teacher's privileged context.
Our diagnosis (Section~\ref{sec:diagnosis}) showed that COT references produce the strongest epistemic suppression (80\% reduction), making this the ideal setting to test whether masking can prevent degradation.

\begin{table*}[t]
\centering
\small
\begin{tabular}{lcccc}
\toprule
\textbf{Method} & \textbf{Step} & \textbf{AIME24} & \textbf{AIME25} & \textbf{MATH-500} \\
\midrule
Base Model & 0 & 0.272 & 0.208 & 0.696 \\
\midrule
Naive OPSD & 100 & 0.217 & 0.192 & 0.760 \\
Masked OPSD & 100 & 0.250 & 0.208 & 0.740 \\
\midrule
\multicolumn{2}{l}{Naive $\Delta$ (step 0 $\to$ 100)} & \textbf{$-$0.055 ($-$20.2\%)} & $-$0.016 ($-$7.7\%) & $+$0.064 (+9.2\%) \\
\multicolumn{2}{l}{Masked $\Delta$ (step 0 $\to$ 100)} & \textbf{$-$0.022 ($-$8.1\%)} & 0.000 (0\%) & $+$0.044 (+6.3\%) \\
\bottomrule
\end{tabular}
\caption{COT reference solution experiments (DeepSeek-R1-Distill-Qwen-1.5B, 100 steps). Naive OPSD degrades AIME24 by 20.2\%; masked OPSD reduces degradation to 8.1\% and completely preserves AIME25.}
\label{tab:cot_results}
\end{table*}

Table~\ref{tab:cot_results} reports the results.
Naive OPSD with COT references causes a substantial drop in AIME24 accuracy from 0.272 to 0.217 ($-20.2\%$), confirming that rich reference context can indeed degrade reasoning.
AIME25 also declines (0.208 $\to$ 0.192, $-7.7\%$).
In contrast, \textbf{masked OPSD significantly mitigates the degradation}: AIME24 drops only to 0.250 ($-8.1\%$), and AIME25 remains completely stable at 0.208 (0\% change).
MATH-500, an in-distribution benchmark, improves under both conditions.

The full training trajectories are shown in Figure~\ref{fig:cot_aime}.
[INSERT FIGURE: AIME24 and AIME25 curves for Naive vs Masked COT OPSD across 100 steps]

\subsection{Clean Reference Solution Experiments (7B)}

We also evaluate on the original OPSD setup with concise reference solutions on DeepSeek-R1-Distill-Qwen-7B.
Our diagnosis (Table~\ref{tab:ref_context}) showed that clean solutions produce only mild epistemic suppression, suggesting that degradation in this setting should be less severe.

[INSERT TABLE: Clean-solution results on 7B — Naive OPSD AIME24 0.478 $\to$ 0.481 (flat) vs.\ Masked OPSD similarly stable]

The results confirm that with concise reference solutions, naive OPSD on 7B does not produce the dramatic degradation observed with COT references or with SDPO's regeneration mechanism.
This is consistent with our diagnostic finding that degradation severity depends on the richness of the reference context: concise solutions suppress epistemic tokens by only 37\%, while full COT traces suppress by 80\%.
Nevertheless, masking provides additional stability and, more importantly, the diagnostic framework we provide can predict \textit{when} masking will be necessary.

\subsection{Analysis}

\paragraph{Epistemic token retention.}
We track the frequency of epistemic tokens in student responses throughout training.
[INSERT PLOT: epistemic token count per response over steps]

\paragraph{Response length dynamics.}
Naive OPSD with COT references tends to reduce response length as the model learns to imitate the teacher's concise style.
Masked OPSD maintains healthier response lengths.
[INSERT PLOT]

\paragraph{Per-token KL divergence.}
Under masked OPSD, the per-token KL divergence at non-masked positions remains stable at 0.03--0.05 throughout training, indicating that masking a small fraction of tokens does not degrade the distillation signal at informative positions.


% ═══════════════════════════════════════════════════════════════
\section{Discussion}
% ═══════════════════════════════════════════════════════════════

\paragraph{Why does such a simple intervention work?}
Our diagnostic findings provide a clear causal account: self-distillation degradation is driven by a concentrated mechanism—the teacher's systematic suppression of a small, specific set of tokens.
These epistemic markers constitute only 1--2\% of all response tokens but account for a disproportionate share of the KL divergence.
By protecting only these positions, we eliminate the source of degradation without meaningfully reducing the distillation signal elsewhere.

\paragraph{Domain generality.}
Our analysis reveals why self-distillation succeeds in code generation and scientific tool use while failing in mathematical reasoning.
Code tasks rely on syntactic correctness and pattern matching, where epistemic uncertainty markers are neither present nor needed.
In such domains, the teacher's confident style is actually beneficial—it eliminates unnecessary hesitation.
However, in mathematical reasoning, where exploration, self-verification, and recovery from intermediate errors depend on explicit uncertainty expression, suppressing epistemic tokens is catastrophic.
Our masking method provides a domain-adaptive fix: it preserves epistemic expression when it matters (reasoning) and does not interfere when it does not (code).

\paragraph{Reference context as a control knob.}
Our experiments reveal that the type of privileged information provided to the teacher modulates both the degradation severity and the mask's effectiveness.
Concise solutions produce mild suppression (37\%); full COT traces cause extreme suppression (80\%).
This suggests that practitioners can choose reference contexts to balance learning speed and reasoning preservation, and that masking becomes increasingly important as the reference context becomes richer.

\paragraph{Limitations.}
Our study is limited to the DeepSeek-R1-Distill-Qwen model family.
While our experiments and prior work suggest that the epistemic suppression effect generalizes across diverse model families, we have not yet validated masking on architectures beyond DeepSeek-R1-Distill-Qwen.
The predefined epistemic token set may also miss novel uncertainty expressions, though our experiments suggest the current set captures the most impactful tokens.
Adaptive token identification and per-domain epistemic sets are promising directions for future work.

\paragraph{Broader implications.}
Our findings reinforce a key insight: optimizing solely for answer correctness can silently reshape reasoning behavior in harmful ways.
Even when the training objective is mathematically sound, the learned behavior may overfit to the distribution of privileged information available at training time.
Simple behavioral safeguards—such as protecting uncertainty tokens—can prevent such pathologies without complex architectural changes.
We hope our diagnostic methodology and principled intervention encourage closer examination of \textit{how} post-training objectives reshape reasoning, not just \textit{whether} they improve accuracy.


% ═══════════════════════════════════════════════════════════════
\section{Conclusion}
% ═══════════════════════════════════════════════════════════════

We presented a systematic investigation into why self-distillation degrades mathematical reasoning and a simple, effective solution.
Through token-level probability probes, we demonstrated that the teacher's privileged context causes systematic suppression of epistemic tokens—the uncertainty markers essential for complex reasoning.
This suppression intensity varies with reference context type, explaining the domain split between self-distillation's success in code and its failure in math.
Our proposed fix, epistemic token masking, requires no additional hyperparameters, no architectural changes, and no extra compute, making it a drop-in addition to any self-distillation pipeline.
Experiments across two model scales and two reference-solution types confirm that masking significantly reduces AIME24 degradation (from $-20\%$ to $-8\%$) and completely preserves AIME25 performance.
We hope our work encourages the community to pay closer attention to how post-training objectives reshape reasoning behavior—and to the power of simple, targeted interventions on small sets of key tokens.


% ═══════════════════════════════════════════════════════════════
% Bibliography
% ═══════════════════════════════════════════════════════════════

\bibliography{custom}
\bibliographystyle{acl_natbib}

\end{document}
